\documentclass[xcolor=dvipsnames]{beamer}

\definecolor{commentgreen}{rgb}{0,0.6,0}
\definecolor{cordovagrey}{rgb}{0.2,0.2,0.2}
\definecolor{cordovablue}{rgb}{0.3,0.76,0.89}


\usetheme{Berlin}
\setbeamertemplate{headline} %hide subsection bar at the top
{
	\begin{beamercolorbox}[colsep=1.5pt]{upper separation line head}
	\end{beamercolorbox}
	\begin{beamercolorbox}{section in head/foot}
		\vskip2pt\insertnavigation{\paperwidth}\vskip2pt
	\end{beamercolorbox}%
	\begin{beamercolorbox}[colsep=1.5pt]{lower separation line head}
	\end{beamercolorbox}
}

\usefonttheme{professionalfonts}
\usecolortheme[named=cordovablue]{structure}

\setbeamercolor{itemize item}{fg=cordovablue}
\setbeamercolor{itemize subitem}{fg=cordovablue}
\setbeamercolor{enumerate item}{fg=cordovablue}
\setbeamercolor{enumerate subitem}{fg=cordovablue}

\setbeamertemplate{itemize item}[square]
\setbeamertemplate{itemize subitem}[square]
\setbeamertemplate{navigation symbols}{}%remove navigation symbols
\expandafter\def\expandafter\insertshorttitle\expandafter{\insertshorttitle\hfill\insertframenumber\,} %show frame numbers

\usepackage[utf8]{inputenc}
\usepackage[ngerman]{babel} %german language
\usepackage{graphicx} %insert pictures
%\usepackage{svg} %insert vector graphics
\usepackage[export]{adjustbox} %for columns?
\usepackage{listings} %insert code
\lstset{ %configure code listings
	basicstyle=\ttfamily,
	keepspaces=true, 
	numbers=none,
	commentstyle=\color{commentgreen},
	keywordstyle=\color{blue},
	escapeinside={\%*}{*)}
}
\usepackage{tikz} %diagrams
\usepackage{verbatim} %diagrams
\usepackage{wasysym} %emojis
\usepackage{hyperref} %links

\title[Entwicklung von Hybrid-Apps am Beispiel von Apache Cordova]{Entwicklung von Hybrid-Apps am Beispiel von Apache Cordova}
%\subtitle{}
\author{Patrick Fehling}
\institute{Hochschule für Technik und Wirtschaft Berlin}
\date{\today}

\begin{document}
\maketitle
\frame{\tableofcontents}


\section{Allgemeines}

\begin{frame}\frametitle{Allgemeine Daten}
	\begin{columns}[t,onlytextwidth]
		\column{.7\textwidth}
		\begin{itemize}
			\item mobile application development framework
			\item Anfängliche Entwicklung durch Nitobi
			\item 2011 Übernahme durch Adobe
			\begin{itemize}
				\item "`PhoneGap"'
				\item "`Apache Cordova"'
			\end{itemize}
			\item Basis für ähnliche Projekte z.B. Ionic, Monaca, Intel XDK
			\item Unterstützung durch: Adobe, BlackBerry, IBM, Intel, Microsoft, Mozilla
			\item Ziel-OSs: iOS, Android, Blackberry etc.
		\end{itemize}
		\column{.3\textwidth}
		\centering
		
		\includegraphics[width=0.6\textwidth]{pictures/cordova_logo}
		
		\begin{tikzpicture} 
		\node [inner sep=0pt] at (0,0) {\includegraphics[width=2.0cm]{pictures/Pacifica}};
		\draw [white, rounded corners=0.4cm, line width=0.35cm] 
			(current bounding box.north west) -- 
			(current bounding box.north east) --
			(current bounding box.south east) --
			(current bounding box.south west) -- cycle;
		\end{tikzpicture}
		\begin{tikzpicture}
		\node [inner sep=0pt] at (0,0) {\includegraphics[width=2.0cm]{pictures/CleverBaby}};
		\draw [white, rounded corners=0.4cm, line width=0.35cm] 
			(current bounding box.north west) -- 
			(current bounding box.north east) --
			(current bounding box.south east) --
			(current bounding box.south west) -- cycle;
		\end{tikzpicture}
	\end{columns}
\end{frame}

\section{Installation \& Build}
\begin{frame}[fragile]\frametitle{Installation - Cordova}
Vorraussetzungen:
\begin{itemize}
	\item Node.js
	\item (optional) Entwicklungsumgebungen für Plattformen
\end{itemize}
\begin{lstlisting}[language=Python]
$ npm install -g cordova
$ cordova create MyApp; cd MyApp
$ cordova platform add <platform>
$ cordova run <platform>
#platform: browser, android, ios...

$ cordova requirements
\end{lstlisting}

\end{frame}

\begin{frame}\frametitle{Installation - Plattformen}
	\begin{columns}[t,onlytextwidth]
		\column{.6\textwidth}
			Android:
			\begin{itemize}
				\item Java SDK (JAVA\_HOME setzen)
				\item Android SDK (ANDROID\_HOME setzen)
				\item gradle (z.B. über Android Studio)
			\end{itemize}
			installieren und der "`PATH"'-Umgebungsvariable hinzufügen\newline
			
			iOS:
			\begin{itemize}
				\item auf OSX: \href{https://cordova.apache.org/docs/en/latest/guide/platforms/ios/index.html}{\textcolor{cordovablue}{iOS Platform Guide}}
				\item auf Windows, Linux: {\LARGE \frownie} ?
			\end{itemize}
		\column{.4\textwidth}
		\centering
		\includegraphics[width=0.9\textwidth,valign=t]{pictures/cordova_device_is_ready}
	\end{columns}
\end{frame}

\section{"'Hello World"'-Beispiel}
\section{Grenzen von Hybrid-Apps}
\section{Sensorik-Nutzung}

\begin{frame}
	\centering
	\textcolor{cordovablue}{{\LARGE Vielen Dank für Ihre Aufmerksamkeit!\\[6ex] Fragen?}}
\end{frame}

\begin{frame}\frametitle{Quellen}
\begin{itemize}
	\item https://cordova.apache.org/
	\item https://ccoenraets.github.io/cordova-tutorial/index.html
	\item https://www.thinkpacifica.com/
	\item https://www.mycleverbaby.com/features.html
	\item https://en.wikipedia.org/wiki/Apache\_Cordova
	\item https://developer.ibm.com/dwblog/2017/what-is-apache-cordova-cross-platform-mobile-application-development/
\end{itemize}
\end{frame}

\end{document}